\section{Exact analytic enclosures and point counterexamples}
\label{sec:enclosures}

\status{Implemented using rational arithmetic. The enclosure proposer and checker share the same numerical kernel; a separate Decimal computation supplies finite sanity checks, not a proof of that kernel}

\subsection{A small exponential enclosure with a proved tail}

A ladder obstruction says nothing about the sign of the original function. To produce a genuine counterexample, the worker must exhibit a point in the original domain and prove that the original expression is negative there. The prototype does this with rational enclosures of exponentials.

For $0\le u\le1/2$ and $n\ge0$, define
\[
 S_n(u)=\sum_{j=0}^n\frac{u^j}{j!},\qquad
 a_{n+1}=\frac{u^{n+1}}{(n+1)!},\qquad
 \rho=\frac{u}{n+2}.
\]

\begin{lemma}[Geometric tail enclosure]
\label{lem:exp-bound}
For the above parameters,
\begin{equation}
 S_n(u)\le e^u\le S_n(u)+\frac{a_{n+1}}{1-\rho}.
 \label{eq:exp-bound}
\end{equation}
\end{lemma}
\begin{proof}
The exponential series has nonnegative terms for $u\ge0$, so its partial sum is a lower bound. For $j\ge n+1$, consecutive omitted terms have ratio $u/(j+1)\le\rho<1$. Thus the tail is at most the geometric series $a_{n+1}\sum_{k\ge0}\rho^k=a_{n+1}/(1-\rho)$.
\end{proof}

For arbitrary rational $t$, choose $k\ge0$ such that $u=|t|/2^k\le1/2$. Compute \eqref{eq:exp-bound}, square the positive interval $k$ times, and, if $t<0$, take the reciprocal with reversed endpoints. This uses only rational arithmetic and the elementary identities
\[
 e^{2v}=(e^v)^2,\qquad e^{-v}=1/e^v.
\]
The prototype allows degrees $0\le n\le64$, $|t|\le64$, and an input-denominator bit-length cap of 128. Unsupported inputs fail the numerical worker; these limits are not assumptions in the mathematical enclosure lemma.

There are two important advantages over printing a decimal approximation. Every endpoint is an exact rational number, and the tail bound has a short mathematical justification. There is also a cost: exact squaring can create large numerators and denominators. A future implementation may use proved outward-rounded dyadics for compactness, but that would be a different numerical kernel requiring its own correctness proof.

\subsection{Interval evaluation of the original normal form}

Given a rational interval $[a,b]$, each monomial term $cx^je^{\lambda x}$ is enclosed by interval multiplication of an enclosure for $x^j$ and one for $e^{\lambda x}$. The exponential is monotone, so the two rate-scaled endpoints suffice. For even $j$ and an interval crossing zero, the lower bound on $x^j$ is zero; otherwise the appropriate endpoint powers are used. Multiplication considers all four endpoint products. A reciprocal is permitted only when the input interval excludes zero.

Summing term enclosures yields a valid but potentially wide enclosure of \eqref{eq:exp-poly}. This is a natural interval extension, not a Taylor model of the full expression. It does not preserve correlations between repeated occurrences of $x$, and it can be inconclusive even when the target is far from changing sign. Subdivision or a better normal form may help; changing a failed lower bound into zero would not.

The same limitations explain why a polynomial square worker is still valuable. Expanding $(e^x-2)^2$ and applying interval arithmetic loses the obvious square structure. The existing symbolic worker can establish nonnegativity without evaluating an exponential at all.

\subsection{A point certificate has a different logical meaning}

The point search tries rational candidates $0,1/4,2/4,\ldots,8$, and increases the Taylor degree through $8,16,32$. These are proposal heuristics. A point certificate contains the selected rational point, the degree, and the resulting rational enclosure. Its checker:
\begin{enumerate}
\item decodes the original externally supplied positivity problem;
\item checks that the point lies on its right ray;
\item recomputes the enclosure from the original expression; and
\item accepts only if the upper endpoint is strictly negative.
\end{enumerate}

An accepted point certificate mathematically refutes the universal nonnegativity statement. It is not merely failure of a proof language. Conversely, failure to find such a point within the list and degree budgets says only unknown. In particular, the eight true controls in \eqref{eq:outside} return a ladder obstruction and no point counterexample. The two outcomes must never be conflated.

In an eventual Lean interface, a point counterexample could justify a theorem of the form $\neg(\forall x\ge0,\ 0\le f(x))$, once the enclosure evaluator and its source bridge are verified. A tactic trying to prove $f(x)\ge0$ cannot use that negative theorem to close the positive goal. It can report a diagnosis or discharge a separately requested negation.

\section{Root certificates as analytic witness plans}
\label{sec:roots}

\subsection{The target includes the interval}

The root lane receives an expression $f\in\E$ and an \emph{externally specified} rational interval $[a,b]$, with $a<b$. Its intended conclusion is
\begin{equation}
 \exists! x\in[a,b],\quad f(x)=0,
 \label{eq:unique-root}
\end{equation}
where uniqueness is restricted to that interval. The notation means that there is exactly one real number satisfying both interval membership and the equation.

The producer may not replace a requested global uniqueness theorem by a convenient local interval and claim the original goal has been solved. For a goal asking merely for a root in a larger domain, a smaller bracket is useful only after the host proves that the bracket is contained in the requested domain. This source-to-goal condition is separate from numerical validity.

\subsection{Endpoint signs plus an exact derivative cover}

A certificate chooses a direction $d\in\{-1,1\}$ and proves
\begin{equation}
 df(a)<0<df(b),\qquad df'(x)>0\quad(x\in[a,b]).
 \label{eq:root-signs}
\end{equation}
Strict endpoint signs are certified by rational enclosures. The derivative is recomputed from the original normal form, then enclosed on a finite subdivision tree. Each leaf carries a Taylor degree and a rational interval with strictly positive lower endpoint after multiplication by $d$.

Coverage is checked structurally. The root node inherits exactly $[a,b]$. A split at rational $m$ must satisfy $a<m<b$, and its children inherit $[a,m]$ and $[m,b]$. Leaves do not get to choose their own smaller domains. The checker limits the total tree to 2,048 visited nodes; the proposer uses midpoint splits with maximum depth eight. These checks prevent omitted gaps and detached successful subintervals from being mistaken for a full derivative proof.

\begin{theorem}[Unique-root certificate]
\label{thm:root}
For a continuous function on $[a,b]$, differentiable there, conditions \eqref{eq:root-signs} imply \eqref{eq:unique-root}, with the root in fact lying in $(a,b)$.
\end{theorem}
\begin{proof}
Apply the intermediate value theorem to $df$, whose endpoint values have opposite signs. The root cannot equal either endpoint because their signs are strict. The derivative condition and the mean value theorem make $df$ strictly increasing on the interval. Two distinct roots would contradict strict monotonicity.
\end{proof}

This theorem makes an interval certificate into a useful existence proof. It does not need an exact symbolic closed form for the root. The Mathlib intermediate-value and derivative-monotonicity APIs provide concrete formalization targets~\cite{mathlib-ivt,mathlib-mvt}.

\subsection{Concrete tested root problems}

The corpus includes eight equations
\[
 xe^x=q,\qquad q\in\{1/8,1/4,1/2,1,2,3,4,8\},
\]
with distinct supplied rational brackets, four equations $e^x=q$ with $q=2,3,5,10$, and four additional brackets. Two of the latter isolate two distinct positive roots of $e^x=3x$ separately:
\[
 [1/2,1],\qquad [3/2,2].
\]
The derivative certificate is decreasing on the first bracket and increasing on the second. This pair is a deliberate check against promoting interval uniqueness to global uniqueness.

The remaining two examples are $e^x=2x+1$ on $[1,3/2]$ and $e^{-x}=x$ on $[1/2,3/5]$. For $xe^x=1$, the recorded bracket is $[1/2,3/5]$. The implementation never requires a Lambert-$W$ library: the defining equation, continuity, endpoint enclosures, and derivative sign suffice.

\subsection{An exact existence proof is not a computable real value}

The root certificate proves that an exact real number exists in a bracket. It does not identify the midpoint, a decimal expansion, or either endpoint as that root. A Lean proof of an existential proposition can be constructed by the intermediate value theorem. A computational function returning arbitrary-precision approximations would need an additional representation of computable reals, nested brackets, a convergence argument, and an error specification. None of those is delivered by a single bracket.

For a synthesis system, there are therefore two distinct interfaces. A proposition-valued hole such as
\[
 \exists x\in[1/2,3/5],\ xe^x=1
\]
can be filled by an existence theorem. A data-valued request for a machine number whose exact real interpretation solves the equation cannot be filled by pretending that the bracket midpoint is exact. A noncomputable choice of the unique real is possible in a sufficiently classical Lean context, but its computational and axiom implications must be explicit. Flow does not change the user's permitted axiom policy.

\subsection{Deliberate incomplete cases}

The current root worker requires strict endpoint signs and a strict derivative-sign cover. It returns unknown for an endpoint root, an even-multiplicity interior root with no sign change, or a simple root whose derivative cannot be bounded away from zero by the chosen interval extension. A function may have a unique root even when these sufficient conditions fail. Increasing numerical budgets is sometimes useful, but it is not a completeness theorem.

The guard tests include an endpoint-root case, a forged split, an altered external bracket, and a reversed monotonicity direction. These test different failure modes. They do not amount to a proof that the checker is hardened against every malformed or adversarial byte sequence.
